\cotefile{1}
\scanpage{93}

\uline{Th.~4} (\uline{troisième th. de convexité dans les espaces de Hilbert})

Soient $E,F$ deux espaces de Hilbert, $u,v$ deux opérateurs compacts de $E$ dans $F$, soit enfin $\varphi\inserted{(s_1,\ldots s_n)}$ une fonction \deleted{convexe symétrique}\inserted{convexe symétrique croissante \deleted{et \unsure{épist.}}} de $n$ arguments $s_i\geq 0$, \deleted{telle que \unsure{$f$ l'ét.|$f(e^{t_1},$}, \unsure{étant|$e^{t_n})$} soit fonction convexe de $(h)$ $\in R^n$. Alors,}\inserted{Alors}

\[
\varphi(\rho_1(u+v),\ldots,\rho_n(u+v))\leq
\varphi(\rho_1(u)+\rho_1(v),\ldots,\rho_n(u)+\rho_n(v))\tag*{(41)}
\]

En effet, c'est un cas particulier du Th.~1, $1^{\mathrm o}$, \unsure{hypothèse|hyp.} c, compte tenu des formules du N\textsuperscript{o}~1.

\deleted{\uline{Corollaire} \unsure{Soit|Soient} $u$ et $v$ étant comme ci-dessus, soit $\varphi(s)$ une fonction convexe \inserted{croissante} de $s\geq 0$, alors on a pour tout entier $n\geq 0$}

\[
\deleted{\sum_{1}^{n}\varphi(\rho_i(u+v))\leq\sum_{1}^{n}\varphi(\rho_i(u)+\rho_i(v))}\tag*{\deleted{(42)}}
\]

Nous serons surtout \unsure{intéressé|intéressés} par le

\uline{Corollaire} 1 Soient \deleted{$u$ et $v$ comme ci-dessus,}\inserted{$E$ et $F$ deux espaces de Hilbert,} et soit $N$ \uline{une norme de Schatten} sur $R^n$, \deleted{alors la quantité} posons, pour \deleted{tout} opérateur compact $u$ de $E$ dans $F$

\[
\lVert u\rVert_N=N(\rho_1(u),\ldots\rho_n(u))\tag*{(4\deleted{\unsure{2}}3)}
\]

Alors on a

\[
\deleted{\lVert\lambda u\rVert_N=|\lambda|\lVert u\rVert_N}\quad
\boxed{\lVert u+v\rVert_N\leq\lVert u\rVert_N+\lVert v\rVert_N}\tag*{(44)}
\]

\deleted{\unsure{et} $\unsure{u\to\lVert u\rVert_N}$ \unsure{définit|est donc} \illegible{2} $u\in L_0(E,F)$ \unsure{tels que|sous-espace} $N$ est une norme de Schatten \unsure{sur} $\unsure{R_+^{(n)}|R^{(N)}}$}

En effet, $N(s_1,\ldots s_n)$ pour $s_i\geq 0$ satisfait aux conditions de (41), on a donc

\[
N(\rho_1(u+v),\ldots,\rho_{\unsure{i|n}}(u+v))\leq
N(\rho_1(u)+\rho_1(v),\ldots,\rho_{\unsure{i|n}}(u)+\rho_{\unsure{i|n}}(v))
\]

et le second mb est majoré par $N(\rho_1(u),\ldots,\rho_n(u))+N(\rho_1(v),\ldots,\rho_n(v))$ puisque $N$ est une norme, d'où (44).

En particulier

\uline{Corollaire 1} \deleted{\unsure{Pour les|Soient}} si $u,v\in L_0^p(E,F)$, \deleted{\unsure{et si}} (cf $1\leq p\leq+\infty$) \unsure{on a} pour \deleted{\unsure{tout}} \inserted{entier $n>0$}

\[
\deleted{\lVert u+v\rVert_p\leq}
\left(\sum_{1}^{n}(\rho_i(u+v))^p\right)^{1/p}\leq
\left(\sum_{1}^{n}\rho_i(u)^p\right)^{1/p}+
\left(\sum_{1}^{n}\rho_i(v)^p\right)^{1/p}\tag*{(45)}
\]

donc \deleted{\unsure{en passant|par passage} \illegible{3}}

\[
\boxed{\lVert u+v\rVert_p\leq\lVert u\rVert_p+\lVert v\rVert_p}\tag*{(46)}
\]
% archive sheet number: 94
